\title{Parameter-Efficient Orthogonal Finetuning via Factorization}
\begin{equation}
\begin{aligned}
    \!\!\bm{B}(d)=\tilde{\bm{B}}(d,d)\!\cdot\!\begin{bmatrix}
    \bm{B}_1(\frac{d}{2})\!\!\!\!\! & \bm{0} \\
    \bm{0}\!\!\!\!\! &\bm{B}_2(\frac{d}{2})
    \end{bmatrix}= \tilde{\bm{B}}(d,d)\tilde{\bm{B}}(d,\frac{d}{2})\cdots\tilde{\bm{B}}(d,2),
\end{aligned}
\end{equation}

\begin{equation}\label{eq:but_factor}
\begin{aligned}
    \bm{B}^F(k)=\begin{bmatrix}
    \text{diag}(\bm{d}_1) & \text{diag}(\bm{d}_2) \\
    \text{diag}(\bm{d}_3) & \text{diag}(\bm{d}_4)
    \end{bmatrix}\in\mathbb{R}^{k\times k},
\end{aligned}
\end{equation}

\begin{theorem}[Expressivity of BOFT]\label{thm:exp_boft}
    With the same block size, BOFT demonstrates greater expressive power than OFT. By sequentially multiplying butterfly matrices as $\bm{B}_{d-1,1}(d)\bm{B}^\top_{d-1,2}(d)\cdots\bm{B}_{1,1}(d)\bm{B}^\top_{1,2}(d)$, where each $\bm{B}_{i,j}(d), \forall i, \forall j$ is a butterfly matrix, one can approximate any orthogonal matrix of size $d$.
\end{theorem}

\begin{equation}
    \begin{aligned}\nonumber
    \bm{z}=\big(\bm{R}(m,b)\cdot\bm{W}^0\big)^\top\bm{x},~~~~\text{s.t.}~\bigg{\{}\bm{R}(m,b)=\prod_{i=1}^m \tilde{\bm{B}}^b_{(d,i)}~~\&~~\underbrace{\big(\tilde{\bm{B}}^b_{(d,j)}\big)^\top\tilde{\bm{B}}^b_{(d,j)}=\tilde{\bm{B}}^b_{(d,j)}\big(\tilde{\bm{B}}^b_{(d,j)}\big)^\top\!=\bm{I}_d}_{\forall j\in [1, m]}\bigg{\}},
    \end{aligned}
\end{equation}

To tackle this problem, we introduce a method that constructs a dense orthogonal matrix by multiplying several sparse orthogonal matrices together. To systematically analyze the sparsity structure involved in orthogonal matrix factorization, we recast the task of generating a dense orthogonal matrix in OFT as an analogy to an information transmission process. In detail, the construction of a dense matrix $\bm{R}\in\mathbb{R}^{d\times d}$ as a product of $m$ square matrices, written as $\thickmuskip=2mu \medmuskip=2mu \bm{R}=\bm{B}_m\bm{B}_{m-1}\cdots\bm{B}_1$, can be represented by the transfer of information across a lattice of $d\times (m+1)$ nodes, as depicted in Figure~\ref{fig:info_trans}. The rationale for employing this information transmission framework stems from the realization that a dense $d \times d$ square matrix embodies the full interconnectivity from one set of $d$ nodes to another set. For any entry $R_{ij}$ in matrix $\bm{R}$, a value of zero signifies the absence of a pathway for information to pass from node $j$ to node $i$, while a nonzero value allows such transmission. Consequently, expressing the dense matrix $\bm{R}$ as a sequence of matrices $\bm{B}_m\bm{B}_{m-1}\cdots\bm{B}_1$ corresponds to a sequence of information exchanges shaped by the connectivity graphs determined by each $\bm{B}_i$. The information thus flows first through $\bm{B}_1$, sequentially traversing up to $\bm{B}_n$. For a tangible illustration (see Figure~\ref{fig:info_trans}), we consider the specific factorization $\bm{R}=\bm{B}_5\bm{B}_4\bm{B}_3\bm{B}_2\bm{B}_1$, where both the pattern of nonzero entries and the resulting induced graph are displayed. The figure shows the unrolled product of the matrices. In this context, each matrix $\bm{B}_i$ serves as an adjacency matrix linking nodes between the $i$-th and $(i+1)$-th layers; that is, the $(j_1, j_2)$-th entry in $\bm{B}_i$ indicates whether there is a directed connection from node $j_2$ at level $i$ to node $j_1$ at level $(i+1)$ (a zero denotes no such edge). To ensure $\bm{B}_5\bm{B}_4\bm{B}_3\bm{B}_2\bm{B}_1$ forms a dense matrix, it is necessary for each node at the first level to be capable of relaying information to any node at the sixth level. However, if we examine $\bm{R}=\bm{B}_4\bm{B}_3\bm{B}_2\bm{B}_1$—linking source nodes on the first layer and destination nodes on the fifth—we observe, for instance, that information from node $1$ cannot reach node $3$. This reveals that in the product $\bm{B}_4\bm{B}_3\bm{B}_2\bm{B}_1$, the $(3,1)$ entry is necessarily zero. This pattern generalizes: considering shorter products such as $\bm{R}=\bm{B}_3\bm{B}_2\bm{B}_1$ or $\bm{R}=\bm{B}_2\bm{B}_1$, the mapping between the graph induced by the multiplication and the underlying sparsity pattern remains consistent. Broadly speaking, for $\bm{R}\in\mathbb{R}^{d\times d}$ to be dense, the $m$ matrices $\bm{B}_m,\ldots,\bm{B}_1$ must collectively establish directed connections over a $d\times (m+1)$ grid. Each connection links nodes in adjacent layers (i.e., neighboring columns), with the requirement that every source node in the first layer can reach all target nodes in the $(m+1)$-th layer.

\textbf{Contrast with butterfly-based sparse training}. Numerous studies~\cite{dao2019learning,dao2019kaleidoscope,chen2022pixelated,dao2022monarch} have explored sparse training approaches utilizing butterfly parameterizations. The majority of these works focus on directly expressing the neural network weight matrices via butterfly parameterization and then training the networks from the ground up. In contrast, \cite{dao2022monarch} examines the adaptation of pretrained weights by initially projecting them onto a form of butterfly matrices, followed by fine-tuning the resulting components for various downstream tasks. BOFT distinguishes itself by introducing an alternative fine-tuning approach, where weight transformations occur via weight matrices shared across layers.

\section{Orthogonal Parameterization by Butterfly Factorization}

Model finetuning fundamentally aims to maintain similarity between the pretrained and finetuned models according to specified criteria, thus retaining the acquired pretraining knowledge. Conventional finetuning strategies frequently employ a low learning rate, an empirical practice that limits deviations between the pretrained and finetuned models. Given that weight matrices possess an inherent structure, a more systematic perspective involves safeguarding their relational characteristics, specifically the pairwise angles between neurons. This perspective underpins the Orthogonal Finetuning~(OFT)~\cite{qiu2023controlling} methodology, which applies the same orthogonal transformation to all neurons within a layer, thereby rigorously preserving the angular relationships among neurons during finetuning. OFT has achieved favorable results regarding generalization and convergence in the finetuning of text-to-image diffusion models~\cite{qiu2023controlling}. Nonetheless, representing high-dimensional orthogonal matrices leads to a substantial number of trainable parameters in OFT. To mitigate this, OFT adopts a block-diagonal parameterization, which significantly decreases the parameter count. This improvement in parameter efficiency, however, comes with drawbacks: the imposed block-diagonal structure results in a fixed sparsity pattern and orthogonal transformations act independently within each block. Such a structure, though economical in parameter use, can embed unwanted inductive biases (\emph{e.g.}, diminished expressivity of the block-diagonal orthogonal matrix, which limits its ability to approximate general linear transformations). Moreover, determining an effective sparsity configuration for these blocks is still an unresolved challenge.

In this formulation, we use $\tilde{\bm{B}}^b(d,2^{m - i+1})$ and simply denote it as $\tilde{\bm{B}}^b_{(d,i)}$. The matrix $\bm{I}_d$ refers to the identity matrix of dimension $d$. The orthogonal matrix $\bm{R}(m,b) \in \mathbb{R}^{d\times d}$ is constructed as the product of several orthogonal butterfly modules. For notational brevity, BOFT configured with $\bm{R}(m,\frac{b}{2})$ is denoted as BOFT($m$,$b$), with the constraint $b \geq 2$. When $m=1$, BOFT($1$,$b$) coincides with the block-diagonal OFT~\cite{qiu2023controlling} where the blocks are of size $b$. Similarly, BOFT($1$,$d$) becomes the original OFT~\cite{qiu2023controlling} in which the orthogonal matrix has no block restrictions and is fully dense. When taking $m = \log \frac{2d}{b}$ and $b$ arbitrary, the block butterfly matrix $\bm{B}^{\frac{b}{2}}(d)$ acts as $\bm{R}$, creating a dense orthogonal transformation. Broadly, BOFT($m$,$b$) incorporates $\frac{1}{2}(b-1)dm$ learnable parameters when adapting a linear transformation of shape $d\times n$. Employing a butterfly matrix structure (\emph{i.e.}, $m=\log d$, $b=2$), BOFT’s parameter count is reduced to $\mathcal{O}(d\log d)$. In comparison, the traditional OFT with a dense orthogonal matrix involves $\mathcal{O}(d^2)$ parameters, while the block-diagonal OFT, partitioned into $r$ blocks, has $\mathcal{O}(bd)$. Consequently, constructing a dense orthogonal matrix using OFT necessitates choosing $b=d$ for the block size, whereas BOFT achieves this for any valid $b$.

\section{Introduction}

where $\bm{d}_i\in\mathbb{R}^{\frac{k}{2}},\forall i$ are arbitrary vectors. For $d = 2^N$, we then construct the $d$-dimensional \emph{butterfly matrix} $\bm{B}(d)\in\mathbb{R}^{d\times d}$ recursively as:

\section{An Information Transmission View on Orthogonal Finetuning}

\textbf{Butterfly structures}. The radix-2 Cooley-Tukey algorithm decomposes the $N$-point discrete Fourier transform into two subtransforms of size $\frac{N}{2}$, inherently resulting in a butterfly architecture expressible as a product of multiple sparse matrices (collectively termed a butterfly matrix). Such butterfly matrices have been adopted for parameterizing orthogonal matrices—circumventing the need for pivoting in Gaussian elimination and promoting computational efficiency~\cite{parker1995random}—as well as stabilizing training in recurrent neural networks~\cite{jing2017tunable} and improving kernel approximation~\cite{munkhoeva2018quadrature}. Additionally, fast linear transformations have been learned using butterfly parameterizations~\cite{dao2019learning,dao2019kaleidoscope}, and butterfly matrices have enabled efficient neural network training~\cite{chen2022pixelated,dao2022monarch}. Their utility extends further to rapid matrix-vector multiplication~\cite{michielssen1996multilevel,o2010algorithm}, compressive matrix approximation~\cite{li2015butterfly}, and network data transmission~\cite{klasing1994broadcasting,li2003linear,rajasingh2016transmission}. Departing from previous explorations, our objective is to leverage butterfly structures to further increase the parameter efficiency of OFT.

Specifically, OFT learns an orthogonal matrix $\bm{R}\in\mathbb{R}^{d\times d}$ for a pretrained linear transformation $\bm{W}^0\in\mathbb{R}^{d\times n}$, modifying the original forward propagation from $\bm{z}=(\bm{W}^0)^\top\bm{x}$ to $\bm{z}=(\bm{R}\bm{W}^0)^\top\bm{x}$, with $\bm{x}\in\mathbb{R}^{d}$ as the input and $\bm{z}\in\mathbb{R}^{n}$ as the output. Zero initialization is accomplished by setting $\bm{R}$ to the identity matrix at the outset. To maintain the orthogonality of $\bm{R}$ throughout finetuning, we adopt the Cayley parameterization following \cite{qiu2023controlling,liu2021orthogonal}, that is, $\bm{R}=(\bm{I}+\bm{Q})(\bm{I}-\bm{Q})^{-1}$, where $\bm{Q}$ is a skew-symmetric matrix satisfying $\bm{Q}=-\bm{Q}^\top$. For improved parameter efficiency, OFT represents the orthogonal matrix $\bm{R}$ in block-diagonal form, namely $\text{diag}(\bm{R}_1,\bm{R}_2,\cdots,\bm{R}_r)$, where each $\bm{R}_i\in\mathbb{R}^{b\times b}$ is a small orthogonal matrix, and $br=d$. However, this block-diagonal parameterization comes with a caveat—it assumes that each neuron's vector (i.e., a column in $\bm{W}^0$) is divided into $r$ equally sized groups, with each block transforming one group independently via its own orthogonal matrix. Although empirical results suggest block-diagonal structures are effective, grouping a neuron's dimensions strictly by index may not be meaningful, highlighting the potential benefits of using dense orthogonal matrices. This raises an important question: \emph{Is it possible to construct a dense orthogonal matrix while preserving parameter efficiency?}

\textbf{BOFT Identity Initialization}. Typical parameter-efficient finetuning procedures aim to initialize from the exact state of a pretrained model, preserving proximity in function space; an example is LoRA’s zero initialization for the low-rank components. In the case of BOFT, all butterfly factors are initialized as identity matrices (i.e., their corresponding skew-symmetric matrices are set to zero when applying the Cayley transform).

Groundbreaking models such as ChatGPT~\cite{brown2020language,chen2021evaluating} and Stable Diffusion~\cite{rombach2022high} exemplify the extraordinary generalization properties that modern large-scale foundation models exhibit. The rapid strides in these models’ capabilities are accompanied by substantial growth in parameter counts

This work introduces Orthogonal Butterfly, a general, parameter-efficient fine-tuning technique for foundation models founded on the butterfly architecture. The central principle for achieving improved parameter efficiency lies in representing a dense orthogonal matrix as a product of several sparse orthogonal matrices. To systematically obtain suitable matrix decompositions, we present a graphical framework centered on information transmission. Within this formalism, the butterfly configuration emerges as a powerful candidate for producing sparse orthogonal matrix factorizations meeting our requirements. BOFT’s empirical utility is demonstrated across fine-tuning tasks for large language models, expansive vision models, and text-to-image generative models. Our experimental results further establish the general superiority of BOFT as a flexible fine-tuning solution.

\textbf{Expressivity of BOFT}. The butterfly construction combined with permutations can exactly reproduce a variety of well-known efficient linear transforms~\cite{dao2019learning,dao2019kaleidoscope} (\emph{e.g.}, fast Fourier transform, Hadamard transform). However, it is still unclear how closely our orthogonal butterfly matrices can approximate arbitrary orthogonal matrices. To explore this, we carry out a simulation task where we attempt to approximate a random dense orthogonal matrix~\cite{anderson1987generation} of dimension $512\times 512$. Figure~\ref{fig:para_eff} shows averages over 10 random initializations; the y-axis displays the approximation error, while the x-axis indicates the number of effective trainable parameters. Each curve with identical color corresponds to a BOFT configuration with fixed block size. The leftmost point for each color refers to BOFT($1$,$b$) (\emph{i.e.}, a block-diagonal OFT of block size $b$). BOFT consistently achieves improved parameter efficiency compared to OFT. For instance, BOFT($9$,$2$) surpasses BOFT($1$,$16$) in expressive power despite using substantially fewer parameters. Typically, selecting a smaller $b$ with a larger $m$ leads to higher parameter efficiency in BOFT. As another example, BOFT($6$,$4$) achieves similar approximation accuracy to BOFT($2$,$16$) but with a smaller number of parameters. Broadly, the butterfly matrix forms a more structured subset of the orthogonal group compared to the block-diagonal variant, establishing that BOFT is provably more expressive than OFT under identical block size.

Nonetheless, BOFT does have certain limitations. Specifically, since its resulting orthogonal matrix is constructed from the multiplication of multiple orthogonal factors, training involves slightly higher computational overhead compared to OFT. Addressing the training time efficiency of BOFT remains an area for future investigation. On the positive side, after fine-tuning is completed, the orthogonal matrices derived by BOFT can be fused into the pretrained model, incurring no additional inference-time latency. Another open question concerns whether the butterfly network represents the most optimal strategy for information transmission. Our graphical information transmission framework opens avenues to borrow ideas from computer networking—a domain with extensive study of topological efficiency for information dissemination. We foresee that future research could identify more optimal network architectures for assembling dense orthogonal matrices.

\section{Intriguing Insights and Discussions}

In contrast to the block-diagonal design adopted by OFT, which balances model flexibility and structural regularity, BOFT leverages the butterfly structure to enable a more gradual transition between matrices in the full orthogonal group (for expressivity) and the identity matrix (for regularity). This design choice permits the selection of a more compact hypothesis space within the orthogonal group when solving downstream problems. Given the prevalence of butterfly architectures in accelerating a range of linear transforms—including the discrete Fourier and cosine transforms—our approach seeks to introduce an intrinsic inductive bias through parameter structure, thus improving generalizability and mitigating overfitting. This hypothesis is rooted in the fact that the butterfly structure readily recovers a variety of canonical linear transforms, unlike the block-diagonal structure used in OFT, which cannot reproduce such transforms. Our key contributions are summarized as follows:

\textbf{Parameter-efficient finetuning (PEFT)}. The growing scale and capability of foundation models (\emph{e.g.}, \cite{brown2020language,radford2021learning,rombach2022high,kirillov2023segment}) have led to significant interest in parameter-efficient techniques for adapting these models to new tasks~\cite{treviso2023efficient,ding2023parameter,lialin2023scaling}. Within the diverse landscape of PEFT approaches~\cite{houlsby2019parameter,aghajanyan2020intrinsic,hulora2022,edalati2022krona,wang2022adamix,gheini2021cross,zaken2022bitfit,guo2020parameter,sung2021training,ansell2022composable,lester2021power,li2021prefix,vu2022spot,he2021towards,mao2021unipelt,karimi2021compacter,liu2022few,sung2022lst,chen2022parameter,jia2022visual,chen2022adaptformer,zhang2022neural,jie2023fact,lian2022scaling,luo2023towards}, reparameterization-based techniques~\cite{aghajanyan2020intrinsic,hulora2022,edalati2022krona} are particularly pertinent to our research. LoRA~\cite{hulora2022}, for instance, modifies pretrained weights by incorporating a product of two low-rank matrices, yielding competitive results in language-related tasks. Because LoRA employs a fixed rank across all layers, subsequent methods \cite{zhang2022adaptive,valipour2022dylora,zhang2023increlora} have sought to allocate the parameter budget more flexibly by varying the rank layerwise. This straightforward approach of low-rank reparameterization has consequently seen broad adoption~\cite{dettmers2023qlora,zi2023delta,chavan2023one}. Motivated by the role of hyperspherical energy in describing generalization properties~\cite{liu2018learning,liu2021orthogonal}, \cite{qiu2023controlling} introduce orthogonal finetuning as a different and effective strategy for finetuning text-to-image diffusion models. In OFT, an orthogonal transformation is learned at the layer level, resulting in superior generalization and more stable training dynamics compared to LoRA. However, OFT typically involves a larger set of trainable parameters relative to LoRA. Enhancing the parameter efficiency of OFT, therefore, represents an important challenge. Moreover, it remains unclear if OFT can be broadly applied to diverse adaptation tasks outside its original application to text-to-image diffusion models~\cite{qiu2023controlling}. BOFT addresses this by integrating butterfly factorization to increase the parameter efficiency of OFT. With this advance, we are now positioned to showcase the effectiveness of orthogonal finetuning for a wide array of adaptation problems.

\textbf{Spectral properties}. Orthogonal finetuning typically achieves superior spectral characteristics compared to LoRA, as it ensures the spectral norm of the original pretrained weight matrix $\bm{W}^0$ remains intact. To understand this, consider the singular value decomposition: $\bm{W}^0=\bm{U}\bm{\Sigma}\bm{V}^\top$, where $\bm{U}$ and $\bm{V}$ are orthogonal matrices and $\bm{\Sigma}$ contains the singular values along its diagonal. Both OFT and BOFT act by left-multiplying an orthogonal matrix $\bm{R}$, resulting in the finetuned form $\bm{R}\bm{U}\bm{\Sigma}\bm{V}^\top$. This transformation leaves the dominant singular value, namely the spectral norm of $\bm{W}^0$, unchanged. The preservation of the spectral norm has been demonstrated to significantly enhance both training stability and the ability to generalize~\cite{miyato2018spectral,yoshida2017spectral}. Additional mathematical insights are provided in Appendix~\ref{app:sn_property}.

\textbf{Orthogonal finetuning as learning bilinear similarity}. The formulation of BOFT can be expressed as learning a bilinear similarity of the form $\bm{w}^0_i\bm{R}\bm{x}$, where $\bm{w}^0_i$ denotes the $i$-th column vector (neuron) in the weight matrix $\bm{W}^0$. In this perspective, BOFT corresponds to adjusting the bilinear similarity matrix $\bm{R}$, but with the stringent constraint that $\bm{R}$ must be orthogonal. This links BOFT to concepts in distance metric learning~\cite{xing2002distance} and the theory of bilinear forms~\cite{roman2005advanced}.

We begin by reviewing foundational aspects of OFT. OFT adapts the pretrained weight matrix by representing the updated weights as the product of a trainable orthogonal matrix and the original, unaltered pretrained matrix. In contrast to LoRA, which performs weight updates through the addition of a low-rank matrix, OFT instead applies a multiplicative update via an orthogonal matrix. LoRA achieves parameter efficiency with its low-rank parameterization, whereas classic OFT employs a block-diagonal orthogonal matrix~\cite{qiu2023controlling}; parameter usage decreases as the diagonal blocks become smaller. Figure~\ref{fig:comp_lora} provides a direct visual comparison of these methods. The core motivation behind using orthogonal transformations for finetuning is to maintain pairwise neuronal angles~\cite{liu2018learning,liu2021orthogonal,qiu2023controlling}, helping to retain much of the semantic information embedded in the pretrained weights.

\input{rebuttal-results}

\textbf{Inductive bias and generalization in BOFT}. The matrix $\bm{R}(m,b)$ within the BOFT framework typically lies within a structured subset of the orthogonal group, imposing restrictions on the class of learnable hypotheses. Consequently, BOFT naturally imparts an inductive bias. We posit that the particular structured bias resulting from butterfly factorization supports generalization due to its alignment with patterns found in established linear transforms~\cite{dao2019kaleidoscope}—such as the discrete Fourier, sine/cosine, and Hadamard transforms. Additionally, the underlying sparse matrix decomposition in BOFT may introduce further implicit inductive biases~\cite{gunasekar2017implicit,li2018algorithmic,liu2019neural}.

(for instance, GPT-3 comprises approximately 175 billion parameters), rendering the process of training such models from scratch an increasingly formidable challenge. Consequently, advancing the field now hinges on methodologies that facilitate adapting existing foundation models to new applications without the need for resource-intensive retraining. This underscores the need for efficient strategies that repurpose foundation models for specific tasks. Task adaptation is generally approached in three main ways: \emph{model finetuning}~\cite{hulora2022,zaken2022bitfit,guo2020parameter,raffel2020exploring,qiu2023controlling,zhang2022adaptive,chavan2023one}, where the core architecture remains intact and selected parameters are updated; \emph{adapter tuning}~\cite{rebuffi2017learning,houlsby2019parameter,pfeiffer2020adapterfusion,lin2020exploring,he2021towards}, which supplements the model with additional tunable modules; and \emph{prompt tuning}~\cite{li2021prefix,lester2021power}, which incorporates optimizable prefix tokens into the input stream. Among these options, model finetuning stands out for its straightforwardness, potent performance, and—crucially—its ability to avoid additional computational delays during inference.

\textbf{Multiplicative Dropout in BOFT}. LoRA~\cite{hulora2022} incorporates a Dropout operation on the low-rank adaptors to regularize training and limit overfitting. Standard Dropout~\cite{srivastava2014dropout} is directly applicable to LoRA, but it is unsuitable for BOFT, since BOFT modifies weights multiplicatively rather than additively. To mitigate overfitting in this context, we introduce a multiplicative Dropout for BOFT: since $\bm{R}(m,b)$ is formed by composing $m$ orthogonal butterfly blocks, each of which may be rearranged as $2b\times 2b$ block-diagonal orthogonal matrices, multiplicative Dropout functions by randomly selecting a proportion $p_1$ of butterfly blocks and a proportion $p_2$ of diagonal sub-blocks within each butterfly, and subsequently substituting these selected portions with identity matrices.

\end{document}

Establishing a straightforward dense linkage between two layers requires $\mathcal{O}(d^2)$ connections, which corresponds to employing a single dense orthogonal matrix. More specifically, this involves $d^2 - d$ edges (for example, when $d=4$, this amounts to $12$ edges). As illustrated in Figure~\ref{fig:info_trans}, one can implement an achievable matrix factorization with only $10$ edges, which is already more parameter-efficient than a fully dense orthogonal matrix. This approach provides a graphical lens through which we can analyze the parameter-efficiency of OFT, facilitating the exploration and design of viable factorizations. Drawing on topologies developed in the Cooley-Tukey algorithm—specifically, butterfly graphs~\cite{cooley1965algorithm}—we note these structures achieve dense interconnectivity between $d$ source and $d$ target nodes with just $\mathcal{O}(d\log d)$ connections. To illustrate, while the structure in Figure~\ref{fig:info_trans} uses $10$ connections, a butterfly network of the same size needs only $8$. We next discuss how employing the butterfly topology can lead to significant parameter savings.

Theorem~\ref{thm:exp_boft} motivates a natural extension for BOFT, namely: the resulting orthogonal matrix is generalized as $\thickmuskip=2mu \medmuskip=2mu \bm{R}^G(m_1,b_1,m_2,b_2,l)=\bm{R}_{l,1}(m_1,b_1)\bm{R}_{l,2}^T(m_2,b_2)\cdots\bm{R}_{1,1}(m_1,b_1)\bm{R}_{1,2}^T(m_2,b_2)$, with $\bm{R}_{i,j}^T(m,b)$ denoting the orthogonal matrix in BOFT. Setting $m_1=m_2=\log d$, $b_1=b_2=2$, and $l=d-1$ allows $\bm{R}^G(m_1,b_1,m_2,b_2,l)$ to span the entire orthogonal group. This style of matrix composition is known as the kaleidoscope hierarchy~\cite{dao2019kaleidoscope}. Nevertheless, we emphasize that higher expressivity does not guarantee better finetuning results; full finetuning, despite being fully expressive, often delivers subpar empirical outcomes. Thus, a balance between expressivity and regularization underpins robust model generalization in finetuning. BOFT expands the search space for finetuning with structural priors, supporting the identification of more favorable trade-offs between expressivity and regularity.

where $\bm{B}_1(\frac{d}{2})$ and $\bm{B}_2(\frac{d}{2})$ denote two butterfly matrices of dimension $\frac{d}{2}$. We introduce the \emph{butterfly component} as $\thickmuskip=2mu \medmuskip=2mu \tilde{\bm{B}}(d,k)=\text{diag}(\bm{B}^F_1(k),\cdots,\bm{B}^F_{d/k}(k))$, which is constructed as a block-diagonal matrix of size $d\times d$, where each block has dimension $k$ and each $\bm{B}^F_i(k)$ is a butterfly factor as defined in Equation~\ref{eq:but_factor}. With these preliminaries, we now discuss parameterizing an orthogonal matrix with the butterfly matrix structure. This is accomplished by ensuring that each block-diagonal factor $\tilde{\bm{B}}(d,k)$ within the butterfly matrix $\bm{B}(d)$ is itself orthogonal. Beginning with the case where $k=2$, we examine the matrix $\tilde{\bm{B}}(d,2)$. Each $2 \times 2$ block can be readily made orthogonal by employing the Cayley transform (or equivalently, a $2$-dimensional rotation), as established in prior works such as \cite{liu2021orthogonal,qiu2023controlling}. Since the sparsity structure of every butterfly component corresponds to a permutation of the pattern in $\tilde{\bm{B}}(d,2)$, every component may be easily parameterized as an orthogonal matrix. This enables a scalable parameterization of orthogonal matrices, leveraging multiple $2\times 2$ orthogonal blocks. Drawing on the generalization in \cite{chen2022pixelated}, we extend the butterfly matrices to define the block butterfly component $\tilde{\bm{B}}^b(d,k)$, wherein each diagonal element $\bm{d}_i$ is replaced by a $b\times b$ matrix. To ensure orthogonality of $\tilde{\bm{B}}^b(d,2)$, we parametrize every $2b\times 2b$ block as an orthogonal matrix. For $k>2$, the corresponding butterfly components $\tilde{\bm{B}}^b(d,k)$ have their nonzero pattern arranged as block-wise permutations of those in $\tilde{\bm{B}}^b(d,2)$, allowing them to also be orthogonalized in a similar manner. Together, these ideas allow us to formulate the forward pass of BOFT as

To facilitate parameter efficiency within this transmission-inspired paradigm, we draw upon the butterfly network structures integral to the Cooley-Tukey fast Fourier transform algorithm~\cite{cooley1965algorithm}, which permit highly efficient information exchange between nodes~\cite{klasing1994broadcasting}. The butterfly configuration inherent in the Cooley-Tukey framework naturally inspires a specific method of sparse matrix factorization, referred to as butterfly factorization. Through this scheme, a dense $d$-dimensional matrix can be realized as a product of $\mathcal{O}(\log d)$ sparse matrices, where each component possesses $\mathcal{O}(d)$ nonzero elements. By guaranteeing the orthogonality of each constituent sparse matrix, our approach attains a highly efficient orthogonal parameterization with merely $\mathcal{O}(d\log d)$ parameters—a drastic reduction relative to the naive approach. Leveraging this parameter-efficient orthogonal scheme, we introduce Orthogonal Butterfly (BOFT), a new method for finetuning that is both parameter- and computation-efficient. Notably, BOFT generalizes OFT, thereby encompassing it as a particular case within a broader orthogonal finetuning framework. The commonality between block-diagonal and butterfly architectures lies in their inherent sparsity, both embedding structured sparsity into orthogonal matrices to curtail the total number of trainable parameters. It is insightful to contrast this paradigm with LoRA’s~\cite{hulora2022} adoption of low-rank structures; both low-rank and sparse matrices exemplify principal classes of structured matrices~\cite{candes2011robust} engineered for parameter-efficient learning.

\section{Concluding Remarks and Limitations}

The perspective of information transmission offers deeper insight into the sparsity characteristics of the factorization matrices in OFT. Figure~\ref{fig:blk_diag} displays the block-diagonal organization of $\bm{R}$ within the original OFT architecture. Although this arrangement reduces the total number of learnable parameters, it does not enable the construction of a fully dense matrix $\bm{R}$. Our objective is to assemble a dense orthogonal matrix from $m$ sparse orthogonal matrices, seeking to minimize the quantity of truly effective trainable parameters. Through the lens of information transmission, two principal criteria emerge: \emph{(i)} \textbf{dense connectivity}, in which each node at the initial stage possesses at least one pathway to each node at the final stage; and \emph{(ii)} \textbf{minimum free edges}, where the overall edge count is kept as small as possible within the limits imposed by orthogonality. Orthogonality imposes intricate constraints on edges joining consecutive layers. Specifically, for every matrix $\bm{B}_i$ to maintain full rank—a prerequisite for orthogonality—there must be $d$ edges, guaranteeing a one-to-one correspondence between the nodes at the $i$-th layer and those at the $(i-1)$-th layer, ensuring no fewer than $d$ edges connect adjacent levels (see, for instance, the value of 4 in Figure~\ref{fig:info_trans}). These $d$ connections are dictated by orthogonality and do not count toward the trainable edges, since they are fixed and non-trainable (for instance, in a $d\times d$ orthogonal matrix with $d$ non-zero entries, each must take the value $\pm 1$). Orthogonal matrices, by virtue of needing fewer parameters compared to general matrices, introduce added interdependence among the edges. Take, for instance, a $2\times 2$ orthogonal matrix: a zero entry at position $(1,1)$ necessitates a corresponding zero at $(2,2)$—so the absence of one edge can lead to the removal of another. Consequently, the orthogonality condition can both introduce and eliminate certain edges, depending on the configuration. The utility of the information transmission framework is most apparent in OFT when visualizing the non-zero structure of the constructed orthogonal matrix, since in this setting only the locations of the non-zero trainable entries in $\bm{R}$ are significant—their actual magnitudes are unimportant.

Originally, the butterfly structure emerged within the Cooley-Tukey algorithm to enable the efficient computation of the fast Fourier transform. Within the Fourier transform context, a modification at a particular frequency can impact the entire spatial domain—an idea resonant with our problem of information transmission, where complete communication from each node at the first layer to every node at the output layer is desired. Moreover, this butterfly design has gained popularity as a network topology in computer science~\cite{solihin2015fundamentals,li2011linear}, allowing for highly efficient communication patterns. Given that $k \geq 2$ and $k$ is a power of $2$, we formally introduce the \emph{butterfly factor} $\bm{B}^F(k)$ as follows:

\begin{itemize}[leftmargin=*,nosep]
    \item We address the challenge of parameter-efficient orthogonal finetuning by formulating a novel information transmission perspective. Under this paradigm, constructing a parameter-efficient dense orthogonal matrix is recast as an information transmission challenge within a grid-structured graph.
    \item Drawing inspiration from the butterfly configurations central to the Cooley-Tukey algorithm, we introduce Orthogonal Butterfly, a method for parameter-efficient orthogonal finetuning formulated within the information transmission framework.
    \item We offer theoretical arguments clarifying how BOFT effectively decreases the count of trainable parameters without compromising expressiveness or training stability. The inherent matrix factorization in BOFT further enables seamless weight interpolation, as depicted in Figure~\ref{fig:boft_interpolation}.
    \item We pioneer the application of orthogonal finetuning to a range of tasks extending beyond controllable text-to-image synthesis~\cite{qiu2023controlling}, demonstrating its versatility as a general finetuning strategy. Specifically, BOFT is evaluated on a spectrum of adaptation benchmarks in both computer vision and natural language processing, consistently outperforming contemporary state-of-the-art approaches and affirming its exceptional parameter efficiency and generalization performance.
\end{itemize}

A central challenge lies in constructing a dense orthogonal matrix in a manner that remains parameter-efficient. While conventional wisdom suggests this task is daunting—since a dense $d$-dimensional orthogonal matrix typically demands $\mathcal{O}(d^2)$ parameters—our methodology innovatively circumvents this by synthesizing a dense orthogonal matrix from the product of multiple structured sparse matrices. The rationale behind this strategy is rooted in the idea that computation time can be traded off against storage needs, effectively reducing the parameter count. Because we encode the orthogonal transformation as the sequential multiplication of sparse matrices, these operations must be carried out at every step during training. Framing this process through the lens of matrix factorization, we recast the creation of a dense orthogonal matrix as a form of information propagation. In this formulation, synthesizing a dense orthogonal matrix via sparse matrix multiplication is analogous to transmitting information across a grid-like graph topology. Adopting this information transmission viewpoint allows us to devise a wide array of sparse matrix factorizations that minimize the trainable parameters, yet retain sufficient flexibility to yield dense matrix representations.


\begin{abstract}
The widespread adoption of large foundation models has made training them from the ground up increasingly impractical due to their computational demands. Consequently, developing strategies to efficiently tailor these models to specialized tasks is of paramount importance. This work explores Orthogonal Finetuning (OFT), a systematic framework for adapting pretrained models to new tasks. While OFT has proven effective in terms of transferability, its reliance on orthogonal matrices means it still requires a substantial amount of trainable parameters, stemming from the inherently high dimensionality of such matrices. To mitigate this, we analyze OFT through the lens of information transfer, identifying several critical characteristics that support enhanced parameter-efficiency. Drawing inspiration from the Cooley-Tukey algorithm, which underpins efficient information propagation in the fast Fourier transform, we design a more parameter-efficient orthogonal matrix construction employing butterfly structures. Integrating this structure with OFT, we present Orthogonal Butterfly (BOFT), a novel and parameter-efficient approach for model adaptation. Not only does BOFT encapsulate OFT as a particular instance, but it also broadens the orthogonal finetuning landscape. Comprehensive empirical investigations are conducted, showcasing BOFT’s efficacy on adapting large vision transformers, extensive language models, and text-to-image diffusion architectures for a range of visual and linguistic downstream applications.
\end{abstract}